% GENERATED FILE. Edit canonical lessons, then run npm run generate:books.
% canonical-source-hash: fnv1a64:5c65af481e5699d2
% canonical-lessons: GU-C262-ganganvu, GU-C262-pokarvu, GU-C262-ramatgamat, GU-C262-ujani, GU-C262-dubki-marvi

\chapter{A day out with friends}
\label{ch:gu-a-day-out-with-friends}
\hlchaptermodality{262}

\begin{chapteropening}
\textbf{By the end of this chapter:} \emph{I can say how we hummed, called out, swam and had a picnic.}

Complete the last lesson of chapter 262: I can say how we hummed, called out, swam and had a picnic.
\end{chapteropening}

\section[\texorpdfstring{\gu{ગણગણવું} (gaṇgaṇvũ)}{ગણગણવું (gaṇgaṇvũ)}]{\gu{ગણગણવું} (gaṇgaṇvũ) — to hum, to murmur}

\label{lesson:GU-C262-ganganvu}



\begin{quote}
Before the new one: say the Gujarati for garba, then the Gujarati for firecracker.
\end{quote}

\subsection*{You'll want to know: \gu{ગણગણવું}}
\textbf{\gu{ગણગણવું}} — \emph{gaṇgaṇvũ} — \textquotedblleft{}to hum, to murmur\textquotedblright{}.

\begin{quote}
\textbf{\gu{તેણે} \gu{ધીમે} \gu{એક} \gu{ગીત} \gu{ગણગણ્યું}.} — \emph{teṇe dhīme ek gīt gaṇgaṇyũ} — \textquotedblleft{}She quietly hummed a song.\textquotedblright{}
\end{quote}

\begin{cousinweb}
\textbf{\gu{ગણગણવું}} (\emph{gaṇgaṇvũ}) is inherited. It comes from \emph{*ganagana}, \textquotedblleft{}murmur\textquotedblright{}, a reconstructed form that imitates the sound; CDIAL lists Gujarati \emph{gaṇgaṇvũ}, \textquotedblleft{}murmur\textquotedblright{} (CDIAL 4013).
\end{cousinweb}

\begin{grammarlens}[title={What you've built}]
One more word for everyday life at A2.
\end{grammarlens}

\subsection*{Your turn}
\begin{itemize}
\raggedright
  \item \emph{Say it:} \emph{gaṇgaṇvũ}
  \item \emph{Say it:} \emph{gaṇgaṇvũ}, once more
  \item \emph{Say it:} \emph{teṇe dhīme ek gīt gaṇgaṇyũ} — She quietly hummed a song.
  \item \emph{Recall:} say the Gujarati for garba, then the Gujarati for firecracker, then say \emph{gaṇgaṇvũ} again
\end{itemize}

\begin{tcolorbox}[breakable,colback=teal!4,colframe=teal!35!black,title={Before you move on}]
What does \gu{ગણગણવું} mean? (\textquotedblleft{}To hum, to murmur\textquotedblright{}.) Say it once more.
\end{tcolorbox}
\section[\texorpdfstring{\gu{પોકારવું} (pokārvũ)}{પોકારવું (pokārvũ)}]{\gu{પોકારવું} (pokārvũ) — to call out, to shout}

\label{lesson:GU-C262-pokarvu}



\begin{quote}
Before the new one: say the Gujarati for firecracker, then the Gujarati for to hum.
\end{quote}

\subsection*{You'll want to know: \gu{પોકારવું}}
\textbf{\gu{પોકારવું}} — \emph{pokārvũ} — \textquotedblleft{}to call out, to shout\textquotedblright{}.

\begin{quote}
\textbf{\gu{મેં} \gu{મારા} \gu{મિત્રને} \gu{પોકાર્યો}.} — \emph{mẽ mārā mitrane pokāryo} — \textquotedblleft{}I called out to my friend.\textquotedblright{}
\end{quote}

\begin{cousinweb}
\textbf{\gu{પોકારવું}} (\emph{pokārvũ}) is inherited. It comes from \emph{*pukkār}, \textquotedblleft{}shout\textquotedblright{}, a reconstructed form that imitates the sound; CDIAL glosses the Gujarati verb \textquotedblleft{}to shout\textquotedblright{} (CDIAL 8246).
\end{cousinweb}

\begin{grammarlens}[title={What you've built}]
One more word for everyday life at A2.
\end{grammarlens}

\subsection*{Your turn}
\begin{itemize}
\raggedright
  \item \emph{Say it:} \emph{pokārvũ}
  \item \emph{Say it:} \emph{pokārvũ}, once more
  \item \emph{Say it:} \emph{mẽ mārā mitrane pokāryo} — I called out to my friend.
  \item \emph{Recall:} say the Gujarati for firecracker, then the Gujarati for to hum, then say \emph{pokārvũ} again
\end{itemize}

\begin{tcolorbox}[breakable,colback=teal!4,colframe=teal!35!black,title={Before you move on}]
What does \gu{પોકારવું} mean? (\textquotedblleft{}To call out, to shout\textquotedblright{}.) Say it once more.
\end{tcolorbox}
\section[\texorpdfstring{\gu{રમતગમત} (ramatgamat)}{રમતગમત (ramatgamat)}]{\gu{રમતગમત} (ramatgamat) — sports, games}

\label{lesson:GU-C262-ramatgamat}



\begin{quote}
Before the new one: say the Gujarati for to hum, then the Gujarati for to call out.
\end{quote}

\subsection*{You'll want to know: \gu{રમતગમત}}
\textbf{\gu{રમતગમત}} — \emph{ramatgamat} — \textquotedblleft{}sports, games\textquotedblright{}.

\begin{quote}
\textbf{\gu{શાળામાં} \gu{રમતગમતનો} \gu{દિવસ} \gu{હતો}.} — \emph{śāḷāmā̃ ramatgamatno divas hato} — \textquotedblleft{}There was a sports day at school.\textquotedblright{}
\end{quote}

A feminine noun.

\begin{cousinweb}
\textbf{\gu{રમતગમત}} (\emph{ramatgamat}) is built on \emph{ramat}, \textquotedblleft{}game\textquotedblright{}, from the taught \emph{ramvũ}, \textquotedblleft{}to play\textquotedblright{} (Sanskrit \emph{rámatē}, CDIAL 10628), doubled by the rhyming \emph{gamat}, \textquotedblleft{}fun\textquotedblright{}, whose own origin is not settled.
\end{cousinweb}

\begin{grammarlens}[title={What you've built}]
One more word for everyday life at A2.
\end{grammarlens}

\subsection*{Your turn}
\begin{itemize}
\raggedright
  \item \emph{Say it:} \emph{ramatgamat}
  \item \emph{Say it:} \emph{ramatgamat}, once more
  \item \emph{Say it:} \emph{śāḷāmā̃ ramatgamatno divas hato} — There was a sports day at school.
  \item \emph{Recall:} say the Gujarati for to hum, then the Gujarati for to call out, then say \emph{ramatgamat} again
\end{itemize}

\begin{tcolorbox}[breakable,colback=teal!4,colframe=teal!35!black,title={Before you move on}]
What does \gu{રમતગમત} mean? (\textquotedblleft{}Sports, games\textquotedblright{}.) Say it once more.
\end{tcolorbox}
\section[\texorpdfstring{\gu{ઉજાણી} (ujāṇī)}{ઉજાણી (ujāṇī)}]{\gu{ઉજાણી} (ujāṇī) — picnic, outing}

\label{lesson:GU-C262-ujani}



\begin{quote}
Before the new one: say the Gujarati for to call out, then the Gujarati for sports.
\end{quote}

\subsection*{You'll want to know: \gu{ઉજાણી}}
\textbf{\gu{ઉજાણી}} — \emph{ujāṇī} — \textquotedblleft{}picnic, outing\textquotedblright{}.

\begin{quote}
\textbf{\gu{અમે} \gu{નદી} \gu{કિનારે} \gu{ઉજાણી} \gu{કરી}.} — \emph{ame nadī kināre ujāṇī karī} — \textquotedblleft{}We had a picnic on the river bank.\textquotedblright{}
\end{quote}

\begin{cousinweb}
\textbf{\gu{ઉજાણી}} (\emph{ujāṇī}) is a tadbhava. It comes from Sanskrit \emph{udy\'{\={a}}na}, \textquotedblleft{}going out\textquotedblright{} and then \textquotedblleft{}park, garden\textquotedblright{}; CDIAL glosses the Gujarati word \textquotedblleft{}feast\textquotedblright{} (CDIAL 2052).
\end{cousinweb}

\begin{grammarlens}[title={What you've built}]
One more word for everyday life at A2.
\end{grammarlens}

\subsection*{Your turn}
\begin{itemize}
\raggedright
  \item \emph{Say it:} \emph{ujāṇī}
  \item \emph{Say it:} \emph{ujāṇī}, once more
  \item \emph{Say it:} \emph{ame nadī kināre ujāṇī karī} — We had a picnic on the river bank.
  \item \emph{Recall:} say the Gujarati for to call out, then the Gujarati for sports, then say \emph{ujāṇī} again
\end{itemize}

\begin{tcolorbox}[breakable,colback=teal!4,colframe=teal!35!black,title={Before you move on}]
What does \gu{ઉજાણી} mean? (\textquotedblleft{}Picnic, outing\textquotedblright{}.) Say it once more.
\end{tcolorbox}
\section[\texorpdfstring{\gu{ડૂબકી} \gu{મારવી} (ḍūbkī mārvī)}{ડૂબકી મારવી (ḍūbkī mārvī)}]{\gu{ડૂબકી} \gu{મારવી} (ḍūbkī mārvī) — to dive, to take a dip}

\label{lesson:GU-C262-dubki-marvi}



\begin{quote}
Before the new one: say the Gujarati for sports, then the Gujarati for picnic.
\end{quote}

\subsection*{You'll want to know: \gu{ડૂબકી} \gu{મારવી}}
\textbf{\gu{ડૂબકી} \gu{મારવી}} — \emph{ḍūbkī mārvī} — \textquotedblleft{}to dive, to take a dip\textquotedblright{}.

\begin{quote}
\textbf{\gu{અમે} \gu{તળાવમાં} \gu{ડૂબકી} \gu{મારી}.} — \emph{ame taḷāvmā̃ ḍūbkī mārī} — \textquotedblleft{}We took a dip in the lake.\textquotedblright{}
\end{quote}

\begin{cousinweb}
\textbf{\gu{ડૂબકી} \gu{મારવી}} (\emph{ḍūbkī mārvī}) is built from \emph{ḍūbkī}, \textquotedblleft{}a dip\textquotedblright{}, and \emph{mārvī}, \textquotedblleft{}to strike\textquotedblright{}. \emph{Ḍūbkī} is built on \emph{ḍūbvũ}, \textquotedblleft{}to sink\textquotedblright{}, which CDIAL derives from \emph{*ḍubb}, \textquotedblleft{}sink\textquotedblright{}, a reconstructed form (CDIAL 5561).
\end{cousinweb}

\begin{grammarlens}[title={What you've built}]
One more word for everyday life at A2.
\end{grammarlens}

\subsection*{Your turn}
\begin{itemize}
\raggedright
  \item \emph{Say it:} \emph{ḍūbkī mārvī}
  \item \emph{Say it:} \emph{ḍūbkī mārvī}, once more
  \item \emph{Say it:} \emph{ame taḷāvmā̃ ḍūbkī mārī} — We took a dip in the lake.
  \item \emph{Recall:} say the Gujarati for sports, then the Gujarati for picnic, then say \emph{ḍūbkī mārvī} again
\end{itemize}

\begin{tcolorbox}[breakable,colback=teal!4,colframe=teal!35!black,title={Before you move on}]
What does \gu{ડૂબકી} \gu{મારવી} mean? (\textquotedblleft{}To dive, to take a dip\textquotedblright{}.) Say it once more.
\end{tcolorbox}
